
\documentclass[a4paper,11pt]{article}
\usepackage[a4paper, margin=8em]{geometry}

% usa i pacchetti per la scrittura in italiano
\usepackage[french,italian]{babel}
\usepackage[T1]{fontenc}
\usepackage[utf8]{inputenc}
\frenchspacing 

% usa i pacchetti per la formattazione matematica
\usepackage{amsmath, amssymb, amsthm, amsfonts}

% usa altri pacchetti
\usepackage{gensymb}
\usepackage{hyperref}
\usepackage{standalone}

% imposta il titolo
\title{Appunti Comunicazioni Numeriche}
\author{Luca Seggiani}
\date{2026}

% imposta lo stile
% usa helvetica
\usepackage[scaled]{helvet}
% usa palatino
\usepackage{palatino}
% usa un font monospazio guardabile
\usepackage{lmodern}

\renewcommand{\rmdefault}{ppl}
\renewcommand{\sfdefault}{phv}
\renewcommand{\ttdefault}{lmtt}

% disponi il titolo
\makeatletter
\renewcommand{\maketitle} {
	\begin{center} 
		\begin{minipage}[t]{.8\textwidth}
			\textsf{\huge\bfseries \@title} 
		\end{minipage}%
		\begin{minipage}[t]{.2\textwidth}
			\raggedleft \vspace{-1.65em}
			\textsf{\small \@author} \vfill
			\textsf{\small \@date}
		\end{minipage}
		\par
	\end{center}

	\thispagestyle{empty}
	\pagestyle{fancy}
}
\makeatother

% disponi teoremi
\usepackage{tcolorbox}
\newtcolorbox[auto counter, number within=section]{theorem}[2][]{%
	colback=blue!10, 
	colframe=blue!40!black, 
	sharp corners=northwest,
	fonttitle=\sffamily\bfseries, 
	title=~\thetcbcounter: #2, 
	#1
}

% disponi definizioni
\newtcolorbox[auto counter, number within=section]{definition}[2][]{%
	colback=red!10,
	colframe=red!40!black,
	sharp corners=northwest,
	fonttitle=\sffamily\bfseries,
	title=~\thetcbcounter: #2,
	#1
}

% U.D.M
\newcommand{\amp}{\ensuremath{\mathrm{A}}}
\newcommand{\volt}{\ensuremath{\mathrm{V}}}
\newcommand{\meter}{\ensuremath{\mathrm{m}}}
\newcommand{\second}{\ensuremath{\mathrm{s}}}
\newcommand{\farad}{\ensuremath{\mathrm{F}}}
\newcommand{\henry}{\ensuremath{\mathrm{H}}}
\newcommand{\siemens}{\ensuremath{\mathrm{S}}}

% circuiti
\usepackage{circuitikz}
\usetikzlibrary{babel}

% disegni
\usepackage{pgfplots}
\pgfplotsset{width=10cm,compat=1.9}

% disponi codice
\usepackage{listings}
\usepackage[table]{xcolor}

\lstdefinestyle{codestyle}{
		backgroundcolor=\color{black!5}, 
		commentstyle=\color{codegreen},
		keywordstyle=\bfseries\color{magenta},
		numberstyle=\sffamily\tiny\color{black!60},
		stringstyle=\color{green!50!black},
		basicstyle=\ttfamily\footnotesize,
		breakatwhitespace=false,         
		breaklines=true,                 
		captionpos=b,                    
		keepspaces=true,                 
		numbers=left,                    
		numbersep=5pt,                  
		showspaces=false,                
		showstringspaces=false,
		showtabs=false,                  
		tabsize=2
}

\lstdefinestyle{shellstyle}{
		backgroundcolor=\color{black!5}, 
		basicstyle=\ttfamily\footnotesize\color{black}, 
		commentstyle=\color{black}, 
		keywordstyle=\color{black},
		numberstyle=\color{black!5},
		stringstyle=\color{black}, 
		showspaces=false,
		showstringspaces=false, 
		showtabs=false, 
		tabsize=2, 
		numbers=none, 
		breaklines=true
}

\lstdefinelanguage{javascript}{
	keywords={typeof, new, true, false, catch, function, return, null, catch, switch, var, if, in, while, do, else, case, break},
	keywordstyle=\color{blue}\bfseries,
	ndkeywords={class, export, boolean, throw, implements, import, this},
	ndkeywordstyle=\color{darkgray}\bfseries,
	identifierstyle=\color{black},
	sensitive=false,
	comment=[l]{//},
	morecomment=[s]{/*}{*/},
	commentstyle=\color{purple}\ttfamily,
	stringstyle=\color{red}\ttfamily,
	morestring=[b]',
	morestring=[b]"
}

% disponi sezioni
\usepackage{titlesec}

\titleformat{\section}
	{\sffamily\Large\bfseries} 
	{\thesection}{1em}{} 
\titleformat{\subsection}
	{\sffamily\large\bfseries}   
	{\thesubsection}{1em}{} 
\titleformat{\subsubsection}
	{\sffamily\normalsize\bfseries} 
	{\thesubsubsection}{1em}{}

% disponi alberi
\usepackage{forest}

\forestset{
	rectstyle/.style={
		for tree={rectangle,draw,font=\large\sffamily}
	},
	roundstyle/.style={
		for tree={circle,draw,font=\large}
	}
}

% disponi algoritmi
\usepackage{algorithm}
\usepackage{algorithmic}
\makeatletter
\renewcommand{\ALG@name}{Algoritmo}
\makeatother

% disponi numeri di pagina
\usepackage{fancyhdr}
\fancyhf{} 
\fancyfoot[L]{\sffamily{\thepage}}

\makeatletter
\fancyhead[L]{\raisebox{1ex}[0pt][0pt]{\sffamily{\@title \ \@date}}} 
\fancyhead[R]{\raisebox{1ex}[0pt][0pt]{\sffamily{\@author}}}
\makeatother

\begin{document}
% sezione (data)
\section{Lezione del 28-05-26}

% stili pagina
\thispagestyle{empty}
\pagestyle{fancy}

% testo
Abbiamo detto che il segnale in ingresso al demapper, campionato in $k$, è:
$$
z(k) = a_k + n'(k), \quad n'(k) \sim \mathcal{N} \left( 0, \sigma_{n'}^2 \right)
$$
dove $n'(x)$ è il rumore gaussiano al ricevitore, con valori:
$$
\sigma_{n'}^2 = \frac{\sigma_n^2}{\beta g^2(0)}, \quad \sigma_n^2 = \frac{N_0}{2} E_{g_R}
$$

Avevamo detto che riguardo ai filtri di interpolazione ricevitore/trasmettitore valeva l'adattamento:
$$
G_T(f) = G_R(f) = \sqrt{G_{RCR}(f)}
$$
sfruttando il filtro a radice di coseno riazlato. Avevamo quindi calcolato BER e SNR a partire da questo modello. In particolae la BER era propozionale all'area sottesa su una regione di curva gaussiana:
$$
P(e) \sim Q \left( \frac{1}{\sigma_{n'}} \right) = Q \left( \sqrt{\frac{2\beta}{N_0}} \right) = Q \left( \sqrt{\frac{6}{M^2 - 1}{\frac{E_r}{N_0}}} \right)
$$

\subsubsection{Energia sul bit}
Veniamo che l'espressione che abbiamo trovato è per la probabilità di errore sula parola. Per fare lo stesso discorso sul bit \textit{codificato} e non, dobbiamo prendere:
$$
E_r = E_d \log_2(\# M) = r log_2 (\# M) E_b
$$
da cui:
$$
P(e) = Q \left( \sqrt{ \frac{6}{M^2 - 1} log_2(  M) \frac{E_d}{N_0} } \right) = Q \left( \sqrt{ \frac{6 r \log_2(  M)}{M^2 - 1} \frac{E_b}{N_0}} \right)
$$
ricordando che $E_d$ è l'energia del bit codificato, ed $E_b$ è l'energia del bit non codificato.

\subsubsection{Efficienza spettrale}
Abbiamo fatto delle considerazioni alla fine della scorsa lezione sull'efficienza energetica di diversi sistemi (PAM a 2 e 4 bit), e come questo impattava il BER.

Veniamo a definire l'efficienza spettrale $\eta_{SP}$:
$$
\eta_{SP} = \frac{R_b}{B_S}
$$
dove $R_b$ è la frequenza dei bit in ingresso al codificatore di canale, e la $B_S$ è la banda del sistema di comunicazione (dopo il filtro di interpolazione). Questo in qualche modo rappresenta l'utilizzazione della banda necessaria alla comunicazione con i bit veri di informazione.

Quindi, si calcolava a partire della frequenza di bit in ingresso $R_b$, la frequenza dopo il codificatore:
$$
R_d = \frac{R_b}{r}
$$
e infine la frequenza (banda) dopo il mapper:
$$
B_S'= \frac{R_b}{r \log_2 ( M)}
$$
Questa era l'espressione che avevamo dato per la banda in banda base del sistema di comunicazione (prima del filtro).
Aggiunto il filtro a coseno rialzato otteniamo che la banda al ricevitore è:
$$
B_S = \frac{1 + \alpha}{2 T_S} = \frac{1 + \alpha}{2 r \log_2 (M)} \frac{1}{T_b}, \quad T_S = \log_2M \frac{1}{R_d}
$$
da cui l'efficienza spettrale risulta come:
$$
\eta_{SP} = \frac{R_b}{B_S} = \frac{2 r \log_2(M)}{1 + \alpha}
$$

\subsection{Sistemi di comunicazione in banda passante}
Passiamo a discutere, dai sistemi di comunicazione in banda base, i sistemi di comunicazione in \textit{banda passante}, dove il segnale generato dal filtro di trasmissione va in ingresso al modulatore. Questo modulerà (attraverso moltiplicazione per coseno) il segnale ad alta frequenza (solitamente in frequenza radio).

Ricordiamo che il segnale $S_T(t)$ in ingresso al modulatore è pari a:
$$
S_T(t) = \sum_i a_i g_T(t - i T_S)
$$
dove $T_S = \frac{1}{f_s}$ è la solita banda di trasmissione (in banda base), e gli $\{a_k\}$ appartengono ad una PAM ad $M$ livelli (come per la banda base).

Innanzitutto soffermiamoci sulle cose che non cambiano in banda passante: i sistemi in banda passante funzionano in termini di probabilità di errore come gli equivalenti in banda base.

Ciò che cambia sono quindi le considerazioni di banda. Ricordiamo che la banda del segnale in banda passante è il \textit{doppio} della banda del segnale in banda base. Questo perché la banda è definita come l'intervallo dello spettro positivo usato per la comunicazione.

\subsubsection{Considerazioni sulla banda}

Descriviamo la modulazione matematicamente, ricordando che questa equivale a:
$$
S_{RF}(t) = S_T(t) \cos(2 \pi f_0 t)
$$
in dominio tempo, e a:
$$
S_{RF}(f) = \frac{S_T(f - f_0) + S_T(f + f_0)}{4}
$$
in dominio frequenziale (dove il $2$ viene dal fatto che si parla di potenza, non ampiezza). 
La densità spettrale del segnale in banda base, ricordiamo, è:
$$
S_T(f) = \frac{1}{T_S} S_a(f) | G_T(f) |^2 = \frac{1}{T_S} R_a(0) | G_T(f) |^2
$$

Dalle considerazioni sulla banda fatte nella scorsa sezione, col filtro a coseno rialzato la banda passante sfruttata è:
$$
B_{RF} = 2 B_T = \frac{1 + \alpha}{T_S} \ \rightarrow \ \eta_{SP, RF} = \frac{\eta_{SP}}{2} = \frac{r \log_2 (M)}{1 + \alpha} 
$$

\subsubsection{Considerazioni di potenza}
Consideriamo adesso la potenza del segnale trasmesso $S_{RF}(t)$, pari a:
$$
P_{RF} = \int_{-\infty}^{\infty} S_{RF}(f) \, df
$$
da cui tenuto conto che la $S_{RF}$ è 2 volte la $S_T$, diviso $4$, si ha subito:
$$
P_{RF} = \frac{P_S}{2}
$$

\subsubsection{Considerazioni energetiche}
Le considerazioni energetiche seguono lo stesso andamento, cioè dobbiamo ricordare il fattore di $2$:
$$
E_{RF} = P_{RF} T_S = \frac{P_S T_S}{2}
$$

\subsection{Ricevitori di sistemi di comunicazione in banda passante}
Veniamo al ricevitore. Dovremmo inanzitutto prevedere un demodulatore, dato dalla moltiplicazione per:
$$
2 \cos(2 \pi f_0 t)
$$
come ricordiamo dal teorema della modulazione.

Il segnale in ingresso al demodulatore sarà:
$$
r_{RF}(t) = S_{RF}(t) \circledast c(t) + W_{RF}(t)
$$
dove il $c(t)$ è la risposta impulsiva del mezzo e $W_{RF}(t)$ è il rumpore gaussiano bianco alla banda passante, con densità spettrale:
$$
S_{W_{RF}} = \frac{N_0}{2}
$$

Supponiamo che il canale sia non distorcente con $c(t) = \delta(t)$. Allora:
$$
r(t) = 2 r_{RF}(t) \cos(2 \pi f_0t) = 2S_{RF}(t) \cos (2 \pi f_0 t) + W(t)
$$
dove abbiamo definito:
$$
W(t) = 2 W_{RF}(t) \cos(2 \pi f_0 t)
$$

Il rumore, quindi, viene "demodulato" (non che fosse modulato a priori), e una volta demodulato dà la densità spettrale di potenza:
$$
S_W(f) = 4 \cdot \frac{ S_{W_{RF}}(f - f_0) + S_{W_{RF}}(f + f_0) }{4}
$$
dove il $4$ a moltiplicare è dato dalle stesse considerazioni viste per la modulazione (dal $2$ a moltiplicare il termine di demodulazione).

Dal momento che $S_{W_{RF}} = \frac{N_0}{2}$, si ha che $S_W(f)$ dopo il demodulatore è $N_0$ (esattamente da quest'ultima equazione).

Per la componente utile del segnale c'è poco da dire, in quanto le repliche del segnale (date dal teorema della modulazione) vengono filtrate dal filtro di ricezione $g_R(t)$, e quindi non ci preoccupano.

Per quanto riguarda il rumore termico si avrà allora che questo sarà un processo aleatorio $n(t)$, con media nulla e varianza:
$$
\sigma_n = N_0
$$

\subsubsection{Probabilità di errore}
Veniamo alla probabilità di errore. Tenuto conto che:
$$
x(k) = a_k + n(k), \quad n(k) \sim \mathcal{N}(0, N_0)
$$
si ottiene (per una PAM a $M$ livelli):
$$
P(e) \sim Q \left( \sqrt{\frac{6 \frac{E_s}{N_0}}{M^2 - 1}} \right)
$$
che coincide con la probabilità di errore della PAM in banda base.

\end{document}
